
\section{本章内容}

本章介绍本体论的核心组成要素（类、属性、公理、实例）、理论基础（描述逻辑、一阶逻辑）
以及工程领域本体建模的基本概念。

\section{文件说明}

\begin{center}
\begin{tabular}{@{}ll@{}}
\toprule
\textbf{文件} & \textbf{内容} \\
\midrule
\texttt{core-concepts.txt} & 本体核心概念：类、属性、关系示例 \\
\texttt{description-logic.txt} & 描述逻辑公理与约束示例 \\
\texttt{first-order-logic.txt} & 一阶逻辑表达式示例 \\
\texttt{reasoning-examples.txt} & 单调与非单调推理示例 \\
\bottomrule
\end{tabular}
\end{center}

\section{关键概念}

\begin{itemize}
\item \textbf{类 (Class)}: 本体的基础单元，代表领域内的一类事物
\item \textbf{属性 (Property)}: 数据属性（Data Property）与对象属性（Object Property）
\item \textbf{公理 (Axiom)}: 约束与规则，如子类关系（\(\sqsubseteq\)）、等价类（\(\equiv\)）
\item \textbf{实例 (Individual)}: 类的具体实例，即ABox中的数据
\end{itemize}

\section{符号说明}

\begin{center}
\begin{tabular}{@{}ll@{}}
\toprule
\textbf{符号} & \textbf{含义} \\
\midrule
\(\sqsubseteq\) & 子类关系（SubClassOf） \\
\(\equiv\) & 等价类（EquivalentTo） \\
\(\forall\) & 全称量词（Universal Restriction） \\
\(\exists\) & 存在量词（Existential Restriction） \\
\(\land\) & 逻辑与（AND） \\
\(\lor\) & 逻辑或（OR） \\
\(\lnot\) & 逻辑非（NOT） \\
\(\rightarrow\) & 蕴含（Implies） \\
\(\sqcap\) & 交（Intersection） \\
\(\sqcup\) & 并（Union） \\
\bottomrule
\end{tabular}
\end{center}
